\documentclass[11pt]{amsart}
\usepackage[T1]{fontenc}
\usepackage{lmodern}
\usepackage{amsmath,amssymb,amsthm,mathtools}
\usepackage[a4paper,margin=30mm]{geometry}
\usepackage[colorlinks=true,linkcolor=blue,citecolor=blue,urlcolor=blue]{hyperref}

\newtheorem{theorem}{Theorem}[section]
\newtheorem*{theorem*}{Theorem}
\newtheorem{lemma}[theorem]{Lemma}
\newtheorem{corollary}[theorem]{Corollary}
\newtheorem{proposition}[theorem]{Proposition}
\newtheorem*{proposition*}{Proposition}
\newtheorem{remark}[theorem]{Remark}

\theoremstyle{remark}
\newtheorem{fact}[theorem]{Fact}
\newcommand{\stackstag}[1]{\href{https://stacks.math.columbia.edu/tag/#1}{#1}}
\title{$\mathcal{PGF}=\mathcal{GP}$ over commutative Noetherian rings}
\author{Yuhui Hua\textsuperscript{*}}
\thanks{\textsuperscript{*}Corresponding author.}
\address{School of Mathematical Sciences, Shanghai Jiao Tong University, Shanghai 200240, China}
\email{huayuhui@sjtu.edu.cn}
\author{Shun-Jie Li}
\address{School of Mathematical Sciences, Shanghai Jiao Tong University, Shanghai 200240, China}
\email{lishunjie@sjtu.edu.cn}
\author{Chencheng Zhang}
\address{School of Mathematical Sciences, Shanghai Jiao Tong University, Shanghai 200240, China}
\email{zhangchencheng@sjtu.edu.cn}
\subjclass[2020]{Primary 13D05; Secondary 13D02, 13D45}
\keywords{Gorenstein projective modules, projectively coresolved Gorenstein flat modules, totally acyclic complexes, \v{C}ech complexes}
\thanks{The authors were supported by the National Natural Science Foundation of China (No.~12131015).}

\begin{document}
\begin{abstract}
We prove $\mathcal{PGF}=\mathcal{GP}$ over every commutative Noetherian ring.
\end{abstract}
\maketitle

\section{Introduction}

Let $R$ be a ring. 
A module ${}_RM$ is \emph{Gorenstein projective} if it is a cycle of an acyclic projective complex $\mathbf P$ such that $\operatorname{Hom}_R(\mathbf P,Q)$ is acyclic for each projective module ${}_RQ$; 
it is \emph{Gorenstein flat} if it is a cycle of an acyclic flat complex $\mathbf F$ such that $J\otimes_R\mathbf F$ is acyclic for each injective module ${}_RJ$; 
it is \emph{projectively coresolved Gorenstein flat} if it is a cycle of an acyclic projective complex $\mathbf P$ such that $J\otimes_R\mathbf P$ is acyclic for each injective module $J_R$. 
We denote the classes of Gorenstein projective, Gorenstein flat, and projectively coresolved Gorenstein flat modules by $\mathcal{GP}(R)$, $\mathcal{GF}(R)$, and $\mathcal{PGF}(R)$, respectively.

\medskip

Gorenstein projective modules extend Auslander--Bridger's modules of $G$-dimension zero \cite{AB}, see \cite{EJ} for a modern definition.
Enochs, Jenda, and Torrecillas introduced Gorenstein flat modules in 1993 \cite{EJT}. 
It is known that $(\mathcal{GP}(R),\mathcal{GP}(R)^{\perp_1})$ is a hereditary cotorsion pair \cite[Corollary~3.4(1)]{cortes-saroch};
whether $(\mathcal{GP}(R),\mathcal{GP}(R)^{\perp_1})$ is complete, or whether $\mathcal{GP}(R)$ is precovering, remains open in general \cite[Question~7.2]{cox}. 
In 2020, \v{S}aroch and \v{S}\v{t}ov\'{i}\v{c}ek introduced projectively coresolved Gorenstein flat modules,
and proved that $(\mathcal{PGF}(R),\mathcal{PGF}(R)^{\perp_1})$ is a complete hereditary cotorsion pair \cite[Theorem~4.9]{saroch-stovicek}, and thus $\mathcal{PGF}(R)$ is special precovering.
Since $\mathcal{PGF}(R)\subseteq\mathcal{GP}(R)$ and $\mathcal{PGF}(R)$ enjoys good homological properties, it is natural to ask whether $\mathcal{PGF}(R)=\mathcal{GP}(R)$.
By \cite{iacob,saroch-stovicek}, this is equivalent to Holm's earlier question that whether $\mathcal{GP}(R)\subseteq\mathcal{GF}(R)$; see \cite{holm-thesis} (also recorded in \cite[Remark~4.5(4)]{ding-li-mao}).
Recently, a Gorenstein projective module which is not Gorenstein flat, is found over a commutative coherent ring of Krull dimension zero \cite{zhang},
answering Holm's question.

\medskip

Several sufficient conditions on commutative Noetherian rings to satisfy $\mathcal{GP}(R)\subseteq\mathcal{GF}(R)$ has been given.
In 2004, Holm proved $\mathcal{GP}(R)\subseteq\mathcal{GF}(R)$ for right coherent rings of finite big finitistic projective dimension   $\operatorname{FPD}(R)$, see \cite[Proposition~3.1]{holm}. 
In 2011, Murfet and Salarian proved the equivalence of total acyclicity and $F$-total acyclicity for complexes of projective modules over commutative Noetherian rings of finite Krull dimension \cite[Lemma~4.20]{murfet-salarian}.
In 2018, Christensen and Kato proved that totally acyclic complexes of projective modules are $F$-totally acyclic for commutative coherent rings with $\sup\{\operatorname{pd}_R F\mid F\text{ is flat}\}<\infty$; see \cite[Claims~2.5]{christensen-kato}.
They also proved it for commutative Noetherian rings with $\operatorname{pd}_R R_{\mathfrak p}<\infty$ for each $\mathfrak p\in\operatorname{Spec}R$; see \cite[Claims~2.8]{christensen-kato}.
In 2020, Iacob proved that $\mathcal{PGF}(R)=\mathcal{GP}(R)$  on an arbitrary ring $R$ with $\operatorname{fd}_R J<\infty$ for any injective right R-module $J$; see \cite[Proposition~9]{iacob}.
Also in 2020, \v{S}aroch and \v{S}\v{t}ov\'{i}\v{c}ek identified, for a right $\aleph_0$-coherent ring, $\mathcal{PGF}(R)$ with the subclass of Gorenstein projective modules filtered by countably presented Gorenstein projective modules; see \cite[Theorem~4.9]{saroch-stovicek}.

\medskip

In this paper, we show that $\mathcal{PGF}(R)=\mathcal{GP}(R)$ for every commutative Noetherian ring. 
\begin{theorem*}
Let $\mathbf P$ be an acyclic complex of projective modules over a commutative Noetherian ring $R$. Then the following are equivalent.

$(1)$ \ $\mathrm{Hom}_R(\mathbf P, Q)$ is acyclic for each $Q \in \mathrm{Proj}(R)$.

$(2)$ \ $J \otimes_R \mathbf P$ is acyclic for each $J \in \mathrm{Inj}(R)$.
\end{theorem*}

It is known that (2) implies (1) \cite[Lemma~4.20(ii)]{murfet-salarian}. We now prove the converse.

% We prove the equality over every commutative Noetherian ring. Let $\mathbf P$ be a totally
% acyclic complex of projective $R$-modules. For each maximal ideal $\mathfrak m$ of $R$, we
% study the complex
% \[
% \mathbf P_{\widehat{R_{\mathfrak m}}}
% =\widehat{R_{\mathfrak m}}\otimes_R\mathbf P.
% \]
% The \v{C}ech complex on generators of $\mathfrak m$ allows us to carry the Hom-acyclicity test
% from $R^{(\aleph_0)}$ to a bounded \v{C}ech complex over the completion. Matlis duality and
% local cohomology then show that $J\otimes_{\widehat{R_{\mathfrak m}}}\mathbf
% P_{\widehat{R_{\mathfrak m}}}$ is acyclic for each injective $\widehat{R_{\mathfrak
% m}}$-module $J$. Faithfully flat descent gives the corresponding statement over each
% $R_{\mathfrak m}$, and acyclicity over $R$ follows by checking at all maximal ideals.


% For which commutative rings $R$ does $\mathcal{GP}(R)=\mathcal{PGF}(R)$ hold? For commutative coherent rings with Krull dimension 0, the equality can fail \cite{zhang}; for commutative Noetherian rings with finite Krull dimension, the equality holds \cite{holm, murfet-salarian}. We think the condition of Krull dimension is not essential, and can be removed.

\section{Proof of the main theorem}

Fix $\mathbf P$ such that $\mathrm{Hom}_R(\mathbf P, Q)$ is acyclic for each $Q \in \mathrm{Proj}(R)$.


\begin{lemma}\label{lem_Ra}
For $a \in R$, write $R_a = \{1,a,a^2, \ldots\}^{-1} R$. Then $\mathrm{Hom}_R(\mathbf P, R_a)$ is acyclic.
\end{lemma}

\begin{proof}
Apply $\mathrm{Hom}_R(\mathbf P, -)$ to the exact sequence $0 \to R^{(\aleph_0)} \xrightarrow{(e_i -ae_{i+1})} R^{(\aleph_0)} \to R_a \to 0$ and use the fundamental theorem of homological algebra.
\end{proof}

Choose an arbitrary maximal ideal of $R$ together with a set of generators, writing $\mathfrak m = (a_1, a_2,\ldots, a_n)$. Write the \v{C}ech complex as $\check C _\mathfrak m = \bigotimes_{i=1}^n [R \xrightarrow{\text{canonical}} R_{a_i}]$.

\begin{lemma}\label{lem_Cm}
$\mathrm{Hom}_R(\mathbf P, \check C_{\mathfrak m})$ is acyclic.
\end{lemma}

\begin{proof}
Note that $R_a \otimes _R R_b \cong R_{ab}$. By Lemma~\ref{lem_Ra}, $\mathrm{Hom}_R(\mathbf P, (\check C_{\mathfrak m})^{i})$ is acyclic for each $i\in \mathbb{Z}$. Since $\check C_{\mathfrak m}$ is bounded, by induction, $\mathrm{Hom}_R(\mathbf P, \check C_{\mathfrak m})$ is acyclic.
\end{proof}

We now turn to $\widehat{R_\mathfrak m}$, the completion of the local ring $R_\mathfrak m$ at its maximal ideal. Write
\begin{equation*}
(-)_{\widehat{R_\mathfrak m}} : \mathrm{Mod}_{R} \longrightarrow \mathrm{Mod}_{\widehat{R_\mathfrak m}}, \qquad X \longmapsto \widehat{R_\mathfrak m} \otimes _R X.
\end{equation*}

\begin{fact}\label{fact}
If each element of $X \in \mathrm{Mod}_R$ is killed by some power of $\mathfrak m$, then the canonical map $X \to X_{\widehat{R_\mathfrak m}}$ is an isomorphism \cite[Lemme~1]{BeauvilleLaszlo}. The functor $(-)_{\widehat{R_\mathfrak m}}$ is exact, that is, $\widehat{R_\mathfrak m}$ is a flat $R$-module \cite[Proposition~10.14]{AtiyahMacdonald}.
\end{fact}

\begin{lemma}\label{lem_vCm}
$\mathrm{Hom}_{\widehat{R_\mathfrak m}}(\mathbf P_{\widehat{R_\mathfrak m}}, (\check C_{\mathfrak m})_{\widehat{R_\mathfrak m}})$ is acyclic.
\end{lemma}

\begin{proof}
By \cite[Theorem~5.1.20]{brodmann-sharp}, each element of the homology groups of $\check C_{\mathfrak m}$ is killed by some power of $\mathfrak m$, and by Fact~\ref{fact}, the canonical map $\check C_{\mathfrak m} \to (\check C_{\mathfrak m})_{\widehat{R_\mathfrak m}}$ is a quasi-isomorphism. Since $\mathbf{P}$ is a projective complex, the functor $\mathrm{Hom}_R(\mathbf P, -)$ preserves quasi-isomorphisms of bounded complexes. Then by Lemma~\ref{lem_Cm}, $ \mathrm{Hom}_R(\mathbf P, (\check C_{\mathfrak m})_{\widehat{R_\mathfrak m}})$ is acyclic. Finally, we have
\begin{equation*}
\mathrm{Hom}_R(\mathbf P, (\check C_{\mathfrak m})_{\widehat{R_\mathfrak m}}) \cong 
\mathrm{Hom}_{\widehat{R_\mathfrak m}}(\widehat{R_\mathfrak m} \otimes_{R} \mathbf P, (\check C_{\mathfrak m})_{\widehat{R_\mathfrak m}}) \cong 
\mathrm{Hom}_{\widehat{R_\mathfrak m}}(\mathbf P_{\widehat{R_\mathfrak m}}, (\check C_{\mathfrak m})_{\widehat{R_\mathfrak m}}) .\qedhere
\end{equation*}
\end{proof}

In the following discussion, we identify $(\check C_{\mathfrak m})_{\widehat{R_\mathfrak m}}$ with $\check C_{\mathfrak m \widehat{R_\mathfrak m}}$ by the following isomorphisms
\begin{equation*}
(\check C_{\mathfrak m})_{\widehat{R_\mathfrak m}} = \widehat{R_\mathfrak m}\otimes _R \check C_{\mathfrak m}  \cong \bigotimes_{i=1}^n [\widehat{R_\mathfrak m} \xrightarrow{\text{canonical}} (\widehat{R_\mathfrak m})_{a_i}] = \check C_{\mathfrak m \widehat{R_\mathfrak m}}.
\end{equation*}
Here $\widehat{R_\mathfrak m}$ is a local ring with maximal ideal $\mathfrak{m}\widehat{R_\mathfrak m}$, and $\check C_{\mathfrak m \widehat{R_\mathfrak m}}$ is the \v Cech complex induced by $(a_1,a_2,\ldots,a_n)$. Let $\kappa$ be the residue field of $\widehat{R_\mathfrak m}$.
Write $E = E_{\widehat{R_\mathfrak m}}(\kappa)$ and the duality
\begin{equation*}
(-)^\vee : (\mathrm{Mod}_{\widehat{R_\mathfrak m}})^{\mathrm{op}} \to \mathrm{Mod}_{\widehat{R_\mathfrak m}}, \quad X \mapsto \mathrm{Hom}_{\widehat{R_\mathfrak m}}(X, E).
\end{equation*}
Here $E$ is an injective cogenerator \cite[p.~528]{matlis}, so $(-)^\vee$ preserves and reflects exactness.

\begin{lemma}\label{lem_dual_acyclic}
$( \check C_{\mathfrak m \widehat{R_\mathfrak m}})^\vee \otimes_{\widehat{R_\mathfrak m}}\mathbf P_{\widehat{R_\mathfrak m}}$ is acyclic.
\end{lemma}

\begin{proof}
First we show that the canonical map $\check C_{\mathfrak m \widehat{R_\mathfrak m}} \to ( \check C_{\mathfrak m \widehat{R_\mathfrak m}})^{\vee\vee}$ is a quasi-isomorphism. By \cite[Theorem~7.1.3]{brodmann-sharp}, the cohomology groups of $\check C_{\mathfrak m \widehat{R_\mathfrak m}}$ are all Artinian $\widehat{R_\mathfrak m}$ modules. Moreover $\widehat{R_\mathfrak m}$ is complete Noetherian, by \cite[Corollary~4.3(3)]{matlis}, we have $H^i ( \check C_{\mathfrak m \widehat{R_\mathfrak m}}) \cong (H^i( \check C_{\mathfrak m \widehat{R_\mathfrak m}}))^{\vee\vee}$ for each $i\in\mathbb{Z}$. Since $(-)^{\vee}$ is exact, we have $(H^i( \check C_{\mathfrak m \widehat{R_\mathfrak m}}))^{\vee\vee} \cong H^i(( \check C_{\mathfrak m \widehat{R_\mathfrak m}})^{\vee\vee})$ for each $i \in\mathbb{Z}$.
It follows that 
\begin{equation*}
\mathrm{Hom}_{\widehat{R_\mathfrak m}}(\mathbf P_{\widehat{R_\mathfrak m}},  \check C_{\mathfrak m \widehat{R_\mathfrak m}})\longrightarrow
\mathrm{Hom}_{\widehat{R_\mathfrak m}}(\mathbf P_{\widehat{R_\mathfrak m}}, ( \check C_{\mathfrak m \widehat{R_\mathfrak m}})^{\vee\vee})\cong
(( \check C_{\mathfrak m \widehat{R_\mathfrak m}})^\vee \otimes_{\widehat{R_\mathfrak m}} \mathbf P_{\widehat{R_\mathfrak m}})^{\vee}
\end{equation*}
are quasi-isomorphisms. By Lemma~\ref{lem_vCm}, $(( \check C_{\mathfrak m \widehat{R_\mathfrak m}})^\vee \otimes_{\widehat{R_\mathfrak m}} \mathbf P_{\widehat{R_\mathfrak m}})^{\vee}$ is acyclic. Since $(-)^\vee$ reflects exactness, $( \check C_{\mathfrak m \widehat{R_\mathfrak m}})^\vee \otimes_{\widehat{R_\mathfrak m}} \mathbf P_{\widehat{R_\mathfrak m}}$ is acyclic.
\end{proof}

\begin{proposition}\label{prop_key}
$J\otimes_{\widehat{R_\mathfrak m}}\mathbf P_{\widehat{R_\mathfrak m}}$ is acyclic for each injective module $J \in \mathrm{Inj}(\widehat{R_\mathfrak m})$.
\end{proposition}

\begin{proof}
In this proof, \cite[\S~6]{murfet-salarian} serves as the main material. By \cite[(2.4.2)]{lipman}, $(\check C_{\mathfrak m \widehat{R_\mathfrak m}})^\vee$ is a dualizing complex over $\widehat{R_\mathfrak m}$. 
Since $\check C_{\mathfrak m \widehat{R_\mathfrak m}}$ is a bounded complex of flat modules, $(\check C_{\mathfrak m \widehat{R_\mathfrak m}})^\vee$ is a bounded complex of injective modules, hence $(\check C_{\mathfrak m \widehat{R_\mathfrak m}})^\vee$ satisfies the conditions in \cite[\S~6, Setup]{murfet-salarian}. By Lemma~\ref{lem_dual_acyclic}, $(\check C_{\mathfrak m \widehat{R_\mathfrak m}})^\vee \otimes_{\widehat{R_\mathfrak m}}\mathbf P_{\widehat{R_\mathfrak m}}$ is acyclic. 
Then by \cite[Proposition~6.1]{murfet-salarian}, $J\otimes_{\widehat{R_\mathfrak m}}\mathbf P_{\widehat{R_\mathfrak m}}$ is acyclic for each injective $\widehat{R_\mathfrak m}$-module $J$. 
\end{proof}

\begin{remark}
The proposition shows that $\mathcal{PGF}(\widehat{R_\mathfrak m})=\mathcal{GP}(\widehat{R_\mathfrak m})$. 
\end{remark}

\begin{lemma}\label{lem_PRm}
$J\otimes _{R_\mathfrak m} \mathbf P_{R_\mathfrak m}$ is acyclic for each $J \in \mathrm{Inj}(R_\mathfrak m)$.
\end{lemma}

\begin{proof}
The canonical map $J \to \widehat{R_\mathfrak m} \otimes _{R_\mathfrak m} J$ is injective \cite[Theorem~56]{Matsumura}. Take the injective envelope $\widehat{R_\mathfrak m} \otimes _{R_\mathfrak m} J \to K$ over $\widehat{R_\mathfrak m}$. Then $J$ is a direct summand of $K$ (as an $R_\mathfrak m$-module). Thus $J \otimes _{R_\mathfrak m} \mathbf P_{R_\mathfrak m}$ is a direct summand of $K \otimes _{R_\mathfrak m} \mathbf P_{R_\mathfrak m} \cong K \otimes _{\widehat {R_\mathfrak m}} \mathbf P_{\widehat{R_\mathfrak m}}$ (as a complex of $R_\mathfrak m$-modules). Since $K \otimes _{\widehat {R_\mathfrak m}} \mathbf P_{\widehat{R_\mathfrak m}}$ is acyclic by Proposition~\ref{prop_key}, $J \otimes _{R_\mathfrak m} \mathbf P_{R_\mathfrak m}$ is also acyclic.
\end{proof}

\begin{lemma}
$J \otimes_R \mathbf P$ is acyclic for each $J \in \mathrm{Inj}(R)$.
\end{lemma}

\begin{proof}
Since $R$ is Noetherian, $J_\mathfrak m$ is an injective $R_\mathfrak m$-module for each maximal ideal $\mathfrak m$ and $J\in \mathrm{Inj}(R)$, by \cite[\S~18, Lemma~5]{Matsumura}. By Lemma~\ref{lem_PRm}, $(J \otimes_R \mathbf P)_\mathfrak m \cong J_\mathfrak m  \otimes_{R_\mathfrak m } \mathbf P_\mathfrak m$ is acyclic for each maximal ideal $\mathfrak m$, and hence $J \otimes_{R} \mathbf P$ is acyclic.
\end{proof}

\begin{corollary}
    Let $R$ be a commutative Noetherian ring, the following are equivalent for an acyclic complex $\mathbf P$ of projective $R$-modules:
    \begin{enumerate}
        \item $\mathrm{Hom}_R(\mathbf P, Q)$ is acyclic for each $Q \in \mathrm{Proj}(R)$.
        \item $J \otimes_R \mathbf P$ is acyclic for each $J \in \mathrm{Inj}(R)$.
        \item $\mathrm{Hom}_R(\mathbf P, R^{(\aleph_0)})$ is acyclic.
        \item $\mathrm{Hom}_R(\mathbf P,F)$ is acyclic for each $F \in \mathrm{Flat}(R)$.
    \end{enumerate}
\end{corollary}

\begin{proof}
The implications $(4)\Rightarrow(1)\Rightarrow(3)$ are immediate.
Under (3), Lemma~2.1 still applies, and the rest of the proof of the main theorem gives (2). 
Under (2), each $Z^i\mathbf P$ is a $\mathcal{PGF}$
module, and by \cite[Theorem~4.4]{saroch-stovicek},
$\mathrm{Ext}^1_R(Z^i\mathbf P,F)=0$ for each flat module $F$ and each $i\in\mathbb{Z}$.
Applying $\mathrm{Hom}_R(-,F)$ to each $0\to Z^i\mathbf{P} \to \mathbf{P}^i \to Z^{i+1}\mathbf{P}\to 0$ gives (4).
\end{proof}

\begin{thebibliography}{99}

\bibitem{AtiyahMacdonald}
M.~F. Atiyah, I.~G. Macdonald,
Introduction to Commutative Algebra,
Addison-Wesley Ser. Math., Addison-Wesley, Reading, MA, 1969.

\bibitem{AB}
M.~Auslander, M.~Bridger,
Stable Module Theory,
Mem. Amer. Math. Soc., no.~94, Amer. Math. Soc., Providence, RI, 1969.

\bibitem{BeauvilleLaszlo}
A.~Beauville, Y.~Laszlo,
Un lemme de descente,
C. R. Acad. Sci. Paris S\'er. I Math. 320(3)(1995), 335--340.

\bibitem{brodmann-sharp}
M.~P. Brodmann, R.~Y. Sharp,
Local Cohomology: An Algebraic Introduction with Geometric Applications,
2nd ed., Cambridge Stud. Adv. Math., vol.~136,
Cambridge Univ. Press, Cambridge, 2013.

\bibitem{christensen-kato}
L.~W. Christensen, K.~Kato,
Totally acyclic complexes and locally Gorenstein rings,
J. Algebra Appl. 17(3)(2018), 1850039.

\bibitem{cortes-saroch}
M. Cort\'es-Izurdiaga, J. \v{S}aroch,
The cotorsion pair generated by the Gorenstein projective modules
and $\lambda$-pure-injective modules,
Israel J. Math. 273 (2026), 127--156.

\bibitem{cox}
S. Cox,
Approximation theory and elementary submodels,
Bull. Symb. Log., First View (2026), 1--45.

\bibitem{ding-li-mao}
N.~Ding, Y.~Li, L.~Mao,
Strongly Gorenstein flat modules,
J. Aust. Math. Soc. 86(3)(2009), 323--338.

\bibitem{EJT}
E.~E. Enochs, O.~M.~G. Jenda, B.~Torrecillas,
Gorenstein flat modules,
Nanjing Daxue Xuebao Shuxue Bannian Kan 10(1)(1993), 1--9.

\bibitem{EJ}
E.~E. Enochs, O.~M.~G. Jenda,
Gorenstein injective and projective modules,
Math. Z. 220(4)(1995), 611--633.

\bibitem{holm-thesis}
H.~Holm,
Gorenstein projective, injective and flat modules,
M.Sc. thesis, University of Copenhagen, Copenhagen, 2000.

\bibitem{holm}
H.~Holm,
Gorenstein homological dimensions,
J. Pure Appl. Algebra 189(1--3)(2004), 167--193.

\bibitem{iacob}
A.~Iacob,
Projectively coresolved Gorenstein flat and Ding projective modules,
Comm. Algebra 48(7)(2020), 2883--2893.

\bibitem{lipman}
J.~Lipman,
Lectures on local cohomology and duality,
in: Local Cohomology and Its Applications, G.~Lyubeznik (ed.),
Lect. Notes Pure Appl. Math., vol.~226,
Marcel Dekker, New York, 2001, pp.~39--89.

\bibitem{matlis}
E.~Matlis,
Injective modules over Noetherian rings,
Pacific J. Math. 8(3)(1958), 511--528.

\bibitem{Matsumura}
H.~Matsumura,
Commutative Algebra,
2nd ed., Math. Lecture Note Ser., vol.~56,
Benjamin/Cummings, Reading, MA, 1980.

\bibitem{murfet-salarian}
D.~Murfet, S.~Salarian,
Totally acyclic complexes over noetherian schemes,
Adv. Math. 226(2)(2011), 1096--1133.


\bibitem{saroch-stovicek}
J.~\v{S}aroch, J.~\v{S}\v{t}ov\'{i}\v{c}ek,
Singular compactness and definability for $\Sigma$-cotorsion and Gorenstein modules,
Selecta Math. (N.S.) 26(2)(2020), Paper No.~23.

\bibitem{zhang}
C.~Zhang,
Commutative coherent rings with $\mathcal{GP}\not\subseteq\mathcal{GF}$, arXiv:2608.17748v3, 2026.

\end{thebibliography}
\end{document}
